% SOURCE_AUTHORITY: sga5_fr_workpass.tex, lines 11845--11870 (LF-normalized unit).
% SOURCE_SHA256: 791F4EFFC5E02832D5D77ED03518C8156D6F07E4C8238B03545DB93D883FBB28
\begin{statement}{Proposição 4.1}
Suponhamos que a subcategoria $\cC_o$ satisfaça as condições (A), (B), bem como a condição (C) seguinte:
\[
(C)\qquad
\text{Todo diagrama }
\begin{tikzcd}[column sep=large,row sep=large]
A \arrow[r,"u"]&B\\
&L \arrow[u,"v"']
\end{tikzcd}
\]
onde $u$ é um epimorfismo e $L\in\Ob\cC_o$, pode ser completado em um diagrama comutativo
\[
\begin{tikzcd}[column sep=large,row sep=large]
A \arrow[r,"u"]&B\\
L_1 \arrow[u] \arrow[r,"w"']&L \arrow[u,"v"']
\end{tikzcd}
\]
que $L_1\in\Ob\cC_o$ e $w$ é um epimorfismo. Nestas condições, os funtores canônicos
\[
D_o^b(\cC)\longrightarrow D^b(\cC),\qquad
D_o^-(\cC)\longrightarrow D^-(\cC)
\]
são completamente fiéis. Além disso, a imagem essencial de $D_o^b(\cC)$ (resp. $D_o^-(\cC)$) é estável por somandos diretos.
\end{statement}

A demonstração desta proposição utiliza essencialmente um critério de Verdier ([2], 4--2 Teorema) e será dada na primeira exposição do seminário sobre Riemann--Roch (SGA 6 I 2.7).
